% §12.1 From Approximate Sets to Approximate Relations
%
% Monograph register: an approximate relation is an approximate set on the pair
% universe T x U. Model-order note tightened for ch9 consistency (one row per
% pair, collapsing to model order two). Cross-refs ch:approximate-sets,
% sec:from-atom-to-set, def:bernoulli-bool, ch:approximate-functions,
% sec:n-ary-channels, sec:relation-composition, sec:approximate-joins. DEFINED:
% sec:sets-to-relations.

\section{From Approximate Sets to Approximate Relations}
\label{sec:sets-to-relations}

A binary relation $R \subseteq T \times U$ is a set whose elements are ordered
pairs, so the chapter-5 construction (\Cref{ch:approximate-sets}) applies with
the universe $T \times U$ in place of $\Universe$. A relation is its membership
predicate $\SetIndicator{R} : T \times U \to \{\top, \bot\}$
(\Cref{sec:from-atom-to-set}), and an \emph{approximate relation} is the noisy
analogue,
\[
  \tilde R : T \times U \longrightarrow \Bernoulli{\BoolType},
\]
whose evaluation $\tilde R(t, u)$ reports whether $(t, u) \in R$ with the two
rates of \Cref{def:bernoulli-bool},
\[
  \fprate = \Prob{\tilde R(t, u) = \top \Given (t, u) \notin R},
  \qquad
  \fnrate = \Prob{\tilde R(t, u) = \bot \Given (t, u) \in R}.
\]
The block structure is two-block, the in-relation block $R$ and its complement,
one rate per block. The universe grew; the apparatus did not.

Two examples fix the picture. The \emph{friend-of} relation $F \subseteq P
\times P$ on a population $P$ stores known friend pairs in a Bloom-style
structure: a query $\tilde F(\text{Alice}, \text{Bob})$ reads the hashed
positions for the pair, with $\fnrate = 0$ (a stored pair has its bits set) and a
tunable $\fprate$, the chapter-5 Bloom analysis applied to pairs. The
\emph{word-in-document} relation $\tilde R \subseteq W \times D$, the inverted
index of information retrieval, holds $(w, d)$ when word $w$ appears in document
$d$, representable as a Bloom filter per document; the chapter-5
space-versus-accuracy trade carries over verbatim.

The chapter-9 view (\Cref{ch:approximate-functions},
\Cref{sec:n-ary-channels}) is equivalent: $\SetIndicator{R}$ is a binary
function $T \times U \to \BoolType$ whose channel has one row per input pair,
collapsing under Axiom~2 to model order two (in the relation or not), exactly as
chapter~5 collapsed the per-element rates of an approximate set. Whether to read
$\tilde R$ as an approximate set on $T \times U$ or as an approximate binary
function is a matter of which downstream machinery is nearer; both are the same
object. \Cref{sec:relation-composition} adds structure on the universe itself,
the relational operators of join, projection, and selection, with rates
propagating by the chapter-6 composition theorems lifted to pairs;
\Cref{sec:approximate-joins} then recovers the MinHash similarity join of
chapter~10 as one instance.
